\clearpage\begingroup\SourceLeaf{38}\NoSourceHead
\begin{center}
[C.]\\[.8em]
[NACHWEIS DER ERDROTATION.]\\[1.2em]
[1.]\\[.7em]
\textsc{Gerling} an \textsc{Gauss}. Marburg, 31. December 1851.
\end{center}
\noindent .\ .\ .\ .\ .\ .\ .\ .\ .\ .\ .\ .\ .\ .\ .\ .\ .\ .\ .\ .\ .
\par\noindent\hspace*{\parindent}\{Der zweite Gegenstand waren die Pendel-Beobachtungen hinsichtlich der
\textsc{Foucault}schen Beobachtungen[\textsuperscript{*})]\begingroup\renewcommand{\thefootnote}{}\footnotetext{[\textsuperscript{*}) Siehe L. \textsc{Foucault}, \textit{Démonstration physique du mouvement de rotation de la terre au moyen du pendule}, Comptes Rendes, Acad. des Sciences \textbf{32}, Paris 1851, S. 421.]}\endgroup. Die Thatsache an sich war hier an meinem
langen Pendel natürlich gleich constatirt, und ich ärgerte mich nur, dass ich
nicht schon 1842 gleich zwei Visire angebracht hatte; damals, wo ich viele
präparatorische Versuche (mit \textsc{Börsch}, der dann gleich abging) machte, liess
ich das Pendel im Meridian schwingen und beobachtete im ersten Vertical.
Hätte ich zugleich ein Visir im Meridian angebracht, so würde ich die fragliche
Abweichung damals schon bemerkt haben (die ja auch vor 150 Jahren
im Fallen schon in Italien gesehen ist, ohne dass jemand den Grund angab).

Nun habe ich mit \textsc{Knoblauch} zusammen zwar eine Menge von Beobachtungen
gemacht in meinem Verliess (während \textsc{Knobl.} in einer Kirche
noch ein zweites Pendel vorgerichtet, dem Publicum gezeigt und grösstentheils
durch seinen Bedienten hat beobachten lassen). Ich habe aber zu
dieser Art von Beobachtungen, wo man so nahe dabei ist, nie ein rechtes
Zutrauen gehabt, und habe deshalb nichts dagegen, diese Art von Beobachtung
aufzugeben, wozu \textsc{Knobl.} Lust zu haben scheint, weil er seit Ende August
nie mehr mit in den Keller gewollt, und wiederholt geklagt hat, dass das
jetzt Beobachtungen geworden, mit denen sich jedermann schon abgegeben.
\HistoricalFoot{XI1.}{5}
\clearpage\SourceLeaf{39}\HistoricalHead{34}{L}{PHYSIK. BRIEFWECHSEL.}
\noindent Dagegen habe ich eine andere, meine ursprüngliche Beobachtungsweise für
die bessere Jahreszeit, und für den zu hoffenden Fall, dass ich unter meinen
Zuhörern einen tauglichen Helfer finde, oder dass endlich mein 10jähriges
Sollicitiren um einen verpflichteten Gehülfen erhört wird, vorbereitet. — Ich
habe nämlich in dem oberen Zimmer (etwa die Mitte der Länge) zwei Visire
und genaue Maassstäbe an den Wänden angebracht, wodurch also Gestalt
und Lage der Ellipse, sogar noch mit einer Controle für die Güte der Beobachtungen
sich ergeben muss. Dass nämlich die Ellipse hier \textls[120]{regelmässig}
entstehen muss, scheint mir klar schon aus meinen bisherigen Beobachtungen.
Dass sie gegen den Sinn der täglichen Bewegung regelmässig \textls[120]{gehen muss},
scheint mir auch klar. Inwiefern sich aber dies letztere durch Beobachtungen
nachweisen lässt, ist mir noch nicht deutlich. — Ebenso behauptet \textsc{Möbius}
\textls[120]{gesehen} zu haben, dass die Bewegung eine andere gewesen, wenn die erste
Anregung im Meridian gewesen, eine andere, wenn im ersten Vertical angeregt
war. \textsc{Knoblauch} behauptet, immer dasselbe \textls[120]{gesehen} zu haben, also das
Gegentheil von \textsc{Möbius}. — Möglich also, dass Nebenumstände sich dabei der
Haupterscheinung beimengen, und deshalb scheinen mir fernere Beobachtungen,
die aber erst zu Worte kommen, wenn die jetzigen tumultuarischen sich ausgetobt
haben, allerdings noch immer nützlich. — Ich habe m[eines] W[issens]
ziemlich fleissig gelesen, was darüber in den Zeitschriften vorgekommen, und
hat mir noch am meisten Befriedigung gewährt, was \textsc{Coombe} im Philosophical
magazine Juni II 1851[\textsuperscript{*})]\begingroup\renewcommand{\thefootnote}{}\footnotetext{[\textsuperscript{*}) J. A. \textsc{Coombe}, \textit{On the Rotation of the Earth}, Philosoph. Magazine 4. ser. \textbf{1}, London 1851 (Jan.—June), S. 554.]}\endgroup\ sagt; auch \textsc{Dufour}[\textsuperscript{**})]\begingroup\renewcommand{\thefootnote}{}\footnotetext{[\textsuperscript{**}) \textit{Sur les déviations apparentes du plan d'oscillation du pendule dans l'expérience de M. Foucault}; nouvelles expériences faites par M. le général \textsc{Dufour} de concert avec MM. \textsc{Wartmann} et \textsc{Marignac}, Comptes Rendus de l'Académie des Sciences \textbf{33}, Paris 1851, S. 13.]}\endgroup\ scheint zu glauben, dass
solide Beobachtungen noch fehlen. Von \textsc{Secchi}s Beobachtungen habe ich seit
jener Notiz in den Astron. Nachr.[\textsuperscript{***})]\begingroup\renewcommand{\thefootnote}{}\footnotetext{[\textsuperscript{***}) \textit{Auszug aus einem Schreiben des Hrn. Prof. Secchi u.s.w.}, Astronomische Nachrichten \textbf{32}, 1851, Spalte 333, siehe besonders Spalte 335.]}\endgroup\ nichts mehr gelesen. Aber bei allem
scheint mir noch immer eine klare Definition dessen zu fehlen, was bei der
s[o] g[enannten] Pendel-Ebene das eigentlich \textls[120]{bleibende} ist.\}
\par\noindent .\ .\ .\ .\ .\ .\ .\ .\ .\ .\ .\ .\ .\ .\ .\ .\ .\ .\ .\ .\ .
\clearpage\endgroup
